\chapter{Metodología propuesta}
\label{cap:metodologia}

En este capítulo se describe la metodología propuesta para aplicar el TDA a la
clasificación de mamografías, desde la preparación de los datos hasta la
evaluación de los modelos.

\section{Visión general del \textit{pipeline}}
El sistema propuesto consta de cinco etapas:
\begin{enumerate}
  \item Preparación de datos (CBIS-DDSM).
  \item Preprocesado de imágenes y extracción de ROIs.
  \item Extracción de descriptores topológicos.
  \item Clasificación (topológica, \textit{deep learning} y fusión).
  \item Evaluación y visualización.
\end{enumerate}

\begin{figure}[H]
\centering
\resizebox{\textwidth}{!}{%
\begin{tikzpicture}[
    node distance = 0.9cm and 1.1cm,
    stage/.style = {
      draw, rounded corners, align=center, font=\small,
      minimum height=1.1cm, text width=2.2cm, fill=blue!6
    },
    arr/.style = {-Latex, thick}
  ]
  \node[stage] (datos)  {Datos\\[-2pt]\footnotesize CBIS-DDSM};
  \node[stage, right=of datos]  (prep)  {Preprocesado\\[-2pt]\footnotesize ROI, CLAHE};
  \node[stage, right=of prep]   (topo)  {Homología\\persistente};
  \node[stage, right=of topo]   (vec)   {Vectorización\\[-2pt]\footnotesize imagen de persist.};
  \node[stage, right=of vec]    (clf)   {Clasificador\\[-2pt]\footnotesize SVM/RF/CNN};
  \node[stage, right=of clf]    (eval)  {Evaluación\\[-2pt]\footnotesize métricas clínicas};

  \draw[arr] (datos) -- (prep);
  \draw[arr] (prep)  -- (topo);
  \draw[arr] (topo)  -- (vec);
  \draw[arr] (vec)   -- (clf);
  \draw[arr] (clf)   -- (eval);
\end{tikzpicture}%
}
\caption{Esquema general del \textit{pipeline} propuesto: de los datos crudos a
las métricas clínicas, pasando por la extracción de descriptores topológicos.}
\label{fig:pipeline}
\end{figure}

\section{Conjunto de datos}
Se utiliza \mbox{CBIS-DDSM}~\cite{lee2017cbisddsm}. El fichero de metadatos
disponible en el repositorio (\texttt{CBIS-DDSM-All-...\,.xlsx}) contiene la
información de cada estudio: identificador de paciente, tipo de lesión (masa o
calcificación), densidad de la mama (BI-RADS), vista (CC/MLO), y la etiqueta de
patología (\texttt{BENIGN}, \texttt{MALIGNANT}, \texttt{BENIGN\_WITHOUT\_
CALLBACK}).

\begin{itemize}
  \item \textbf{Tarea principal}: clasificación binaria benigno vs.\ maligno.
  \item \textbf{Unidad de análisis}: región de interés (ROI) de la lesión y/o
        imagen completa reescalada.
\end{itemize}

\subsection{Protocolo de evaluación y prevención de fuga de datos}
Se adopta una partición estricta en \textbf{entrenamiento / validación / test}.
Para evitar la fuga de datos (\textit{data leakage}), la separación se realiza
\textbf{a nivel de paciente}: todas las imágenes de un mismo paciente pertenecen
a una única partición. Se emplea estratificación por la etiqueta de patología
para mantener la proporción de clases.

\begin{quote}
\textit{Nota metodológica.} En el material de partida
(\texttt{Demo\_Notebook.ipynb}) se observaba una exactitud de entrenamiento del
100\,\% desde la primera época y una exactitud de validación cercana al 100\,\%,
lo que es un indicio claro de fuga de datos o de un problema en el
\textit{pipeline}. Este trabajo corrige dicho protocolo mediante partición por
paciente, control de la normalización y validación cruzada.
\end{quote}

\section{Preprocesado de imágenes}
El preprocesado incluye:
\begin{itemize}
  \item Conversión y normalización de intensidades.
  \item Recorte de la ROI a partir de las máscaras de segmentación.
  \item Reescalado a un tamaño común (p.\,ej. $224\times224$).
  \item Opcionalmente, realce de contraste (CLAHE) y filtrado de ruido.
\end{itemize}

\section{Extracción de descriptores topológicos}
A partir de cada imagen (o ROI) se calcula la homología persistente mediante una
filtración por \textit{sublevel sets} de la intensidad (para imágenes en escala
de grises) o sobre nubes de puntos derivadas de las microcalcificaciones. Se
obtienen los diagramas de persistencia en dimensiones $0$ y $1$, y a partir de
ellos se generan los vectores de características:
\begin{itemize}
  \item Imágenes de persistencia.
  \item Estadísticos de persistencia (número de puntos, persistencia total,
        máxima, media; curvas de Betti).
\end{itemize}

\section{Modelos de clasificación}
Se comparan tres enfoques:
\begin{enumerate}
  \item \textbf{Topológico}: clasificadores clásicos (SVM, \textit{random
        forest}, regresión logística) sobre los descriptores topológicos.
  \item \textbf{Deep learning}: CNN con \textit{transfer learning} sobre las
        imágenes.
  \item \textbf{Fusión}: combinación de las características topológicas con las
        \textit{embeddings} de la CNN.
\end{enumerate}

\section{Métricas de evaluación}
Dado el carácter clínico del problema, se emplean varias métricas
complementarias:
\begin{itemize}
  \item \textbf{Exactitud} (\textit{accuracy}).
  \item \textbf{Sensibilidad} (\textit{recall}) --- crítica para no pasar por
        alto casos malignos.
  \item \textbf{Especificidad} y \textbf{precisión}.
  \item \textbf{F1-score} y \textbf{AUC-ROC}.
  \item \textbf{Matriz de confusión}.
\end{itemize}
La sensibilidad recibe especial atención, ya que un falso negativo (no detectar
un cáncer) tiene un coste clínico mucho mayor que un falso positivo.
